% 從零開始打造大型語言模型（繁體中文版）
% Build a Large Language Model (From Scratch) - Traditional Chinese Edition
% Original Author: Sebastian Raschka
% Translation assisted by: Claude Code (Anthropic)

\documentclass[12pt,a4paper,openany]{book}

% ===== 中文支援 =====
\usepackage{xeCJK}
\usepackage{fontspec}

% 設定中文字型（需根據系統安裝的字型調整）
\setCJKmainfont{Noto Serif CJK TC}  % 主要字型：思源宋體
\setCJKsansfont{Noto Sans CJK TC}   % 無襯線字型：思源黑體
\setCJKmonofont{Noto Sans Mono CJK TC}  % 等寬字型

% 設定英文字型
\setmainfont{Libertinus Serif}
\setsansfont{Libertinus Sans}
\setmonofont{JetBrains Mono}

% ===== 頁面設定 =====
\usepackage[
    a4paper,
    top=2.5cm,
    bottom=2.5cm,
    left=3cm,
    right=2.5cm,
    headheight=15pt
]{geometry}

% ===== 程式碼高亮 =====
\usepackage{minted}
\usemintedstyle{friendly}

% 程式碼區塊設定
\setminted{
    frame=lines,
    framesep=2mm,
    baselinestretch=1.2,
    bgcolor=lightgray!10,
    fontsize=\small,
    linenos=true,
    breaklines=true,
    breakanywhere=false,
    autogobble=true
}

% 行內程式碼設定
\newmintinline[pyinline]{python}{bgcolor=lightgray!20}
\newmintinline[bashinline]{bash}{bgcolor=lightgray!20}

% ===== 數學支援 =====
\usepackage{amsmath}
\usepackage{amssymb}
\usepackage{amsthm}

% ===== 圖片支援 =====
\usepackage{graphicx}
\usepackage{float}
\usepackage{subfig}
\graphicspath{{./figures/}}

% ===== 表格支援 =====
\usepackage{booktabs}
\usepackage{tabularx}
\usepackage{longtable}
\usepackage{multirow}

% ===== 顏色支援 =====
\usepackage{xcolor}
\definecolor{linkcolor}{RGB}{0,102,204}
\definecolor{citecolor}{RGB}{153,51,153}
\definecolor{urlcolor}{RGB}{0,153,153}

% ===== 超連結支援 =====
\usepackage[
    colorlinks=true,
    linkcolor=linkcolor,
    citecolor=citecolor,
    urlcolor=urlcolor,
    bookmarks=true,
    bookmarksopen=true,
    bookmarksnumbered=true,
    pdfauthor={Sebastian Raschka (原作), Claude Code (翻譯輔助)},
    pdftitle={從零開始打造大型語言模型},
    pdfsubject={Large Language Models, Deep Learning, PyTorch},
    pdfkeywords={LLM, GPT, Transformer, AI, Machine Learning}
]{hyperref}

% ===== 目錄設定 =====
\usepackage{tocloft}
\setcounter{tocdepth}{2}
\setcounter{secnumdepth}{3}

% ===== 頁首頁尾 =====
\usepackage{fancyhdr}
\pagestyle{fancy}
\fancyhf{}
\fancyhead[LE]{\leftmark}
\fancyhead[RO]{\rightmark}
\fancyfoot[C]{\thepage}
\renewcommand{\headrulewidth}{0.4pt}
\renewcommand{\footrulewidth}{0pt}

% ===== 列表設定 =====
\usepackage{enumitem}
\setlist{nosep, leftmargin=*}

% ===== 引用格式 =====
\usepackage[style=apa,backend=biber]{biblatex}
\addbibresource{references.bib}

% ===== 索引支援 =====
\usepackage{makeidx}
\makeindex

% ===== 演算法支援 =====
\usepackage{algorithm}
\usepackage{algpseudocode}

% ===== 程式碼清單 =====
\usepackage{listings}
\lstset{
    basicstyle=\ttfamily\small,
    breaklines=true,
    frame=single,
    backgroundcolor=\color{lightgray!10}
}

% ===== 自訂命令 =====
% 程式碼檔案包含
\newcommand{\includecode}[3][python]{
    \begin{listing}[H]
    \caption{#3}
    \inputminted[linenos,frame=lines,framesep=2mm]{#1}{#2}
    \end{listing}
}

% 重點強調
\newcommand{\important}[1]{\textcolor{red}{\textbf{#1}}}
\newcommand{\note}[1]{\textcolor{blue}{\textit{注意：#1}}}

% 術語標註（首次出現時顯示英文）
\newcommand{\term}[2]{#1（\textit{#2}）}

% ===== 定理環境 =====
\newtheorem{theorem}{定理}[chapter]
\newtheorem{lemma}[theorem]{引理}
\newtheorem{proposition}[theorem]{命題}
\newtheorem{corollary}[theorem]{推論}
\newtheorem{definition}{定義}[chapter]
\newtheorem{example}{範例}[chapter]
\newtheorem{remark}{註記}[chapter]

% ===== 文件資訊 =====
\title{
    {\Huge\bfseries 從零開始打造大型語言模型}\\[1cm]
    {\Large Build a Large Language Model (From Scratch)}\\[0.5cm]
    {\large 繁體中文版（台灣）}
}
\author{
    {\large 原著：Sebastian Raschka}\\[0.3cm]
    {\normalsize 出版社：Manning Publications}\\[0.3cm]
    {\normalsize 翻譯輔助：Claude Code (Anthropic)}
}
\date{\today}

% ===== 文件開始 =====
\begin{document}

% 封面
\frontmatter
\maketitle

% 版權頁
\newpage
\thispagestyle{empty}
\vspace*{\fill}
\begin{center}
\begin{minipage}{0.8\textwidth}
    \textbf{從零開始打造大型語言模型（繁體中文版）}\\[0.5cm]

    \textbf{原作資訊：}\\
    原著：Sebastian Raschka\\
    出版社：Manning Publications\\
    ISBN：978-1633437166\\
    出版日期：2024年9月12日\\
    原始儲存庫：\url{https://github.com/rasbt/LLMs-from-scratch}\\[0.5cm]

    \textbf{翻譯資訊：}\\
    翻譯版本：繁體中文（台灣）\\
    翻譯輔助工具：Claude Code（Anthropic）\\
    翻譯日期：2025年11月\\[0.5cm]

    \textbf{授權：}\\
    原著內容版權歸 Sebastian Raschka 和 Manning Publications 所有\\
    翻譯內容遵循原著的 Apache-2.0 授權\\[0.5cm]

    \textbf{致謝：}\\
    本翻譯專案特別感謝原作者 Sebastian Raschka 提供如此優質的技術書籍，\\
    感謝 Manning Publications 授權和出版原著，\\
    以及 Claude Code（Anthropic）協助進行高品質翻譯。\\[0.5cm]

    \textbf{聲明：}\\
    本翻譯堅持高標準，確保技術準確性和可讀性，\\
    為台灣讀者提供優質的學習資源。\\
    所有技術術語均採用台灣標準用語。
\end{minipage}
\end{center}
\vspace*{\fill}

% 序言
\chapter*{繁體中文版序言}
\addcontentsline{toc}{chapter}{繁體中文版序言}

本書是 Sebastian Raschka 所著《Build a Large Language Model (From Scratch)》的繁體中文（台灣）翻譯版本。本翻譯專案的目標是為台灣讀者提供高品質、準確且易讀的中文版本，協助更多人深入理解大型語言模型的運作原理。

\section*{關於本翻譯}

本翻譯遵循以下原則：

\begin{itemize}
    \item \textbf{準確性優先}：忠實於原文內容，確保技術細節的正確性
    \item \textbf{台灣術語標準}：使用台灣資訊工程界通用的專業術語
    \item \textbf{術語一致性}：維護統一的術語表，確保全書用詞一致
    \item \textbf{可讀性}：採用符合中文閱讀習慣的表達方式
    \item \textbf{完整性}：保留所有原文資訊，包括程式碼、圖表和參考資料
\end{itemize}

\section*{如何使用本書}

本書適合：
\begin{itemize}
    \item 想要深入理解 LLM 運作原理的開發者
    \item 希望從零開始實作 GPT 模型的工程師
    \item 對深度學習和自然語言處理有興趣的研究人員
    \item 想要學習 PyTorch 實作技巧的學習者
\end{itemize}

\section*{程式碼與資源}

所有程式碼範例均可在原作者的 GitHub 儲存庫中找到：
\url{https://github.com/rasbt/LLMs-from-scratch}

本翻譯專案的相關資源：
\begin{itemize}
    \item 術語對照表：\texttt{TERMINOLOGY\_GLOSSARY.zh-TW.md}
    \item 翻譯指南：\texttt{TRANSLATION\_GUIDELINES.zh-TW.md}
\end{itemize}

\vspace{1cm}
\noindent
翻譯團隊\\
2025年11月

% 目錄
\tableofcontents
\listoffigures
\listoftables
\listoflistings

% 主要內容
\mainmatter

% 第 1 章
\chapter{理解大型語言模型}
\label{ch:understanding-llms}

% 章節摘要
\begin{quote}
本章介紹大型語言模型（LLM）的基本概念，包括它們的運作原理、應用場景，以及為什麼要從零開始打造自己的 LLM。
\end{quote}

\section{什麼是大型語言模型？}

\term{大型語言模型}{Large Language Model, LLM} 是一種基於深度學習的人工智慧模型，專門用於理解和生成自然語言文字。這些模型透過在大量文字資料上進行訓練，學習語言的模式、結構和語義。

\subsection{LLM 的關鍵特徵}

大型語言模型具有以下關鍵特徵：

\begin{itemize}
    \item \textbf{規模龐大}：通常包含數億到數千億個參數
    \item \textbf{自監督學習}：從未標記的文字資料中學習
    \item \textbf{情境學習}：能夠根據上下文理解和生成文字
    \item \textbf{少樣本學習}：僅需少量範例即可執行新任務
    \item \textbf{泛用性}：可應用於多種不同的自然語言處理任務
\end{itemize}

\subsection{LLM 的演進歷程}

圖 \ref{fig:llm-timeline} 展示了大型語言模型的發展歷程。

\begin{figure}[H]
    \centering
    % 此處將包含時間軸圖片
    % \includegraphics[width=0.9\textwidth]{ch01/llm-timeline.pdf}
    \caption{大型語言模型的發展時間軸}
    \label{fig:llm-timeline}
\end{figure}

\section{LLM 的運作原理}

\subsection{基本架構}

現代 LLM 大多基於 \term{Transformer}{Transformer} 架構，這是一種於 2017 年提出的深度學習架構。Transformer 的核心是 \term{注意力機制}{Attention Mechanism}，它能讓模型關注輸入序列中的重要部分。

\begin{definition}[Transformer 架構]
Transformer 是一種序列到序列模型架構，主要由以下元件組成：
\begin{itemize}
    \item 自注意力層（Self-Attention Layer）
    \item 前饋神經網路（Feed-Forward Network）
    \item 層正規化（Layer Normalization）
    \item 殘差連接（Residual Connection）
\end{itemize}
\end{definition}

\subsection{訓練過程}

LLM 的訓練通常分為兩個階段：

\begin{enumerate}
    \item \textbf{預訓練（Pre-training）}：
    \begin{itemize}
        \item 在大量未標記文字上訓練
        \item 學習語言的一般模式和知識
        \item 使用自監督學習目標，如下一標記預測
    \end{itemize}

    \item \textbf{微調（Fine-tuning）}：
    \begin{itemize}
        \item 在特定任務的資料上進一步訓練
        \item 調整模型以適應特定應用
        \item 可能包括指令微調或基於人類回饋的強化學習
    \end{itemize}
\end{enumerate}

\section{為什麼要從零開始打造 LLM？}

\subsection{學習目標}

從零開始實作 LLM 有以下好處：

\begin{itemize}
    \item \textbf{深入理解}：透過實作理解每個元件的運作方式
    \item \textbf{靈活性}：能夠根據需求自訂模型架構
    \item \textbf{除錯能力}：更容易診斷和解決問題
    \item \textbf{最佳化}：可針對特定用例進行效能最佳化
\end{itemize}

\begin{example}[教育用 LLM]
本書將引導您建立一個小型但功能完整的 GPT 模型，包含：
\begin{itemize}
    \item 約 124M 參數（與 GPT-2 small 相當）
    \item 可在一般筆記型電腦上訓練
    \item 能夠生成連貫的文字
    \item 可載入預訓練權重進行微調
\end{itemize}
\end{example}

\subsection{實用考量}

雖然訓練大型的生產級 LLM 需要龐大的計算資源，但從零開始打造較小的模型仍然非常實用：

\begin{itemize}
    \item 可用於特定領域的應用
    \item 適合資源受限的環境
    \item 便於實驗和研究新的架構
    \item 有助於理解和評估現有的大型模型
\end{itemize}

\section{本書概覽}

本書將帶領您完成以下學習旅程：

\subsection{章節架構}

\begin{description}
    \item[第 2 章] 處理文字資料 --- 學習如何將文字轉換為模型可處理的格式
    \item[第 3 章] 編寫注意力機制 --- 實作 Transformer 的核心元件
    \item[第 4 章] 從零開始實作 GPT 模型 --- 組合各個元件建立完整模型
    \item[第 5 章] 使用未標記資料進行預訓練 --- 訓練模型學習語言模式
    \item[第 6 章] 針對文字分類進行微調 --- 將模型應用於特定任務
    \item[第 7 章] 針對遵循指令進行微調 --- 訓練模型理解和執行指令
\end{description}

\subsection{學習路徑}

圖 \ref{fig:mental-model} 展示了本書涵蓋的內容架構。

\begin{figure}[H]
    \centering
    % \includegraphics[width=0.8\textwidth]{ch01/mental-model.jpg}
    \caption{LLM 開發的心智模型}
    \label{fig:mental-model}
\end{figure}

\section{先備知識}

\subsection{必要的背景知識}

要充分利用本書，您需要具備：

\begin{itemize}
    \item \textbf{Python 程式設計}：紮實的 Python 基礎
    \item \textbf{基礎數學}：線性代數、微積分和機率論的基本概念
    \item \textbf{（可選）深度學習}：有神經網路的基礎會有幫助
\end{itemize}

\subsection{PyTorch 基礎}

本書使用 PyTorch 作為深度學習框架。如果您不熟悉 PyTorch，附錄 A 提供了快速入門指南。

\begin{remark}
即使您是 PyTorch 新手，也可以跟隨本書學習。所有程式碼都有詳細的解釋，並且逐步建構。
\end{remark}

\section{開發環境設定}

\subsection{硬體需求}

本書的程式碼設計為可在一般筆記型電腦上執行：

\begin{itemize}
    \item \textbf{CPU}：現代多核心處理器
    \item \textbf{RAM}：至少 8GB（建議 16GB）
    \item \textbf{GPU}：可選，但會大幅加速訓練
    \item \textbf{儲存空間}：至少 10GB 可用空間
\end{itemize}

\subsection{軟體需求}

您需要安裝以下軟體：

\begin{itemize}
    \item Python 3.10 或更高版本
    \item PyTorch
    \item 其他 Python 套件（詳見 \texttt{requirements.txt}）
\end{itemize}

詳細的安裝指南請參閱 \texttt{setup} 目錄。

\section{本章小結}

本章介紹了：

\begin{itemize}
    \item 大型語言模型的基本概念和特徵
    \item LLM 的運作原理和訓練過程
    \item 從零開始實作 LLM 的價值和意義
    \item 本書的架構和學習路徑
    \item 必要的先備知識和開發環境設定
\end{itemize}

在下一章中，我們將開始處理文字資料，學習如何將原始文字轉換為模型可以理解的數值表示。

\subsection{延伸閱讀}

\begin{itemize}
    \item Vaswani et al., "Attention is All You Need" (2017)
    \item Radford et al., "Language Models are Unsupervised Multitask Learners" (2019)
    \item Brown et al., "Language Models are Few-Shot Learners" (2020)
\end{itemize}

\subsection{練習題}

本章無程式碼練習。建議：

\begin{enumerate}
    \item 研究並比較不同的大型語言模型（如 GPT、BERT、T5）
    \item 探索 LLM 的應用案例
    \item 設定您的開發環境並確認可以執行 PyTorch 程式碼
\end{enumerate}


% 第 2 章
\include{chapters/ch02}

% 第 3 章
\include{chapters/ch03}

% 第 4 章
\include{chapters/ch04}

% 第 5 章
\include{chapters/ch05}

% 第 6 章
\include{chapters/ch06}

% 第 7 章
\include{chapters/ch07}

% 附錄
\appendix

\include{chapters/appendix-A}
\include{chapters/appendix-B}
\include{chapters/appendix-C}
\include{chapters/appendix-D}
\include{chapters/appendix-E}

% 後記
\backmatter

% 參考文獻
\printbibliography[heading=bibintoc,title={參考文獻}]

% 索引
\printindex

\end{document}
